% s-310-p310-m003; original PDF 310, 311
\Aphorism{s-310-p310-m003}{364}
Solange man jung ist und die Metaphysik noch nicht überwunden hat, empfindet man sich gern als Erwählten heimlicher Lenkung und Lebensbestimmung; das Wirkliche erscheint als das Unerlässliche und zugleich als das Eine grosse Loos unter Millionen Möglichkeiten, unwahrscheinlich wie eins zu unendlich, unmöglich fast, und doch Wirklichkeit, unerlässliche Wirklichkeit, deren Ausbleiben mit Vernichtung gleichbedeutend war. Man zittert wie Jener, der unbewusst durch eine ungeheure Gefahr gelaufen ist; der Nachschrecken schlägt Einen nieder gleich jenem Reiter, der den gefrorenen See überritten hatte. „Wie, wenn du diesen einen Glücksfall verfehlt hättest! und du hättest ihn natürlicherweise verfehlen müssen!“ über diese Logik kommt man nicht hinweg, — dieser Mangel an Gerechtigkeit ist nur als \Sperr{Gnade} zu verstehen, Gnade eines verehrungswürdigen Etwas, das uns an der Wegekreuzung zwischen Heil und Verderben bei der Hand nahm und den Heilsweg führte. Später, kälteren Geistes, glaubt man nicht mehr an diese transcendentale Absichtlichkeit im Schicksale des Einzelnen, auch nicht an die Gefährlichkeit der Krisen und die Einzigkeit der Auswege. Wir haben uns mit einem Wirklichen eingerichtet, es liebgewonnen; möglich, dass wir jetzt an seinem Verluste verbluten würden, gewiss, dass wir durch sein Ausbleiben \Foreign{a priori} nicht zu Grunde gegangen wären. Es war teleologisch nicht unbedingt geboten, dass das Eine kam; wir hätten mit dem Anderen auch weiterleben können und müssen, unter Aneignung eines anderen, vielleicht höheren Ichs. Ich meine, hinter jede Lebenswende und Entscheidung können wir uns zurückversetzen und der Vielfältigkeit und odysseischen Erfindsamkeit des Lebens vertrauen, dass es an dieser Stelle auch andere Wege gab, uns weiterzuführen: jetzt freilich, nachdem wir den einen von ihnen solange gegangen sind, können wir den Querübergang auf einen anderen Weg uns nur als vollständige Durchkreuzung und Verneinung unserer gegenwärtigen Persönlichkeit und Bewegungsrichtung vorstellen.
\par

